Es sei
\mathl{P \in N}{} und
\mathl{P \in U}{} ein Kartengebiet von $M$. Es sei
\maabbdisp {\alpha} { U } {V
} {}
eine Karte mit einer offenen Menge
\mathl{V \subseteq H}{,}
\mathl{H=\R_{\geq 0} \times \R^{n-1}}{.} Diese Karte induziert einen Homöomorphismus
\maabbdisp {} {N \cap U} { \alpha (N \cap U)
} {.}
Diese Karten nehmen wir für $N$. Die Diffeomorphieeigenschaft der Kartenwechsel zu $M$ überträgt sich direkt auf die Karten zu $N$. Es sei
\mathl{P \in N \cap \partial M}{.} Dann wird $P$ unter einer Karte auf einen Randpunkt von $H$ abgebildet, und da die Karten für $N$ Einschränkungen von solchen Karten sind, bleibt diese Eigenschaft erhalten. Wenn umgekehrt
\mathl{P \in \partial N}{} ist, so ist natürlich einerseits
\mathl{P \in N \subseteq M}{.} Andererseits gibt es zu $P$ eine Karte zu $N$, die durch Einschränkung von einer Karte zu $M$ herrührt, in der $P$ auf einen Randpunkt von $H$ abgebildet wird.

 [[Kategorie:Latexseite]]
